\section{From the Model to Real Chess}\label{sec:real}

To test a model value, we keep its tree of pawn moves and replace only the leaves by real results, from the
Lichess tablebase~\cite{LichessTB} up to seven men and beyond from search and our king-and-pawn solver.

\begin{definition}[Real-chess oracle value]\label{def:oracle}
Let $X$ be a pawn component on the ordinary board beside T. An \emph{end position} of its model tree is a
position where a line of play first reaches a passed pawn; its two \emph{sides} are white to move and black to
move. Replace each end position by $+\BIG$ if white wins it in real chess with either side to move, by $-\BIG$
if black does, by $\game{+\BIG}{-\BIG}$ if the side to move wins, and by $0$ if the side to move loses. If a side
is drawn, the \emph{draw treatment} counts that draw as a win for black, as a win for white, or as worth $0$; in
the last case the end position is $\game{\BIG r_W}{\BIG r_B}$, where $r_W,r_B\in\{-1,0,1\}$ are white's results
with white and with black to move. The canonical value of the resulting tree is the \emph{oracle value} $O(X)$.
\end{definition}

\begin{lemma}[Sandwich; \lab{P}]\label{lem:sandwich}
Replace every end position whose real-chess value is not known by $-2\BIG$ (for $O^-$) or by $+2\BIG$ (for
$O^+$), and keep every known one, scored as in Definition~\ref{def:oracle} under the chosen draw treatment. Then
$O^-\le O(X)\le O^+$ whatever the unknown results; if $O^-=O^+$, $O(X)$ is certified.
\end{lemma}
\begin{proof}
Every leaf lies between $-2\BIG$ and $2\BIG$, and raising an option never lowers a
game~\cite{Siegel2013}; induct from the leaves.
\end{proof}

\subsection{The Textbook Candidate in Real Chess}\label{sec:real-cand}

The tree is built under either passed-pawn test: the \emph{textbook tree} or the \emph{e.p.-aware tree}.

\begin{theorem}[Candidate in real chess; \lab{C1}]\label{thm:oracle-cand}
For the candidate \pos{W a2 c2 / B b7 d7} beside T (Fig.~\ref{fig:star2_textbook}):
\begin{enumerate}
\item[(a)] In chess the side to move wins; with black to move only 1\ldots~b5 wins; with white to
  move 1.~c4 wins, 1.~a3, 1.~c3 and all king moves lose, and 1.~a4 is unresolved \lab{C2}.
\item[(b)] In the e.p.-aware tree, for every draw treatment, $O$ has left stop $\ge1$ and right stop $\le-1$, so
  it is not $0$ or $\nim{2}$.
\item[(c)] With draws as wins for black, $O=\game{2\pm1}{-2\pm1}$ in both trees.
\end{enumerate}
\end{theorem}

Real chess thus agrees with $\nim{2}$ on the outcome and on every winning first move resolved so far, but not on
the value: in the e.p.-aware tree, which matches chess, $O$ is hot under every draw treatment, so the candidate's
$\nim{2}$ is a feature of the textbook test, not of chess. In the textbook tree, $\nim{2}$ is excluded
if draws count as wins for either side \lab{C1}, and undecided if draws count $0$. A second candidate, \pos{W a3 b3 c2 / B b5 b7 c6},
$\nim{2}$ under the opponent rule with en passant on, has a positive oracle value under every draw treatment
\lab{C1}; the other non-default $\nim{2}$ positions are open (Section~S7).

\subsection{Towards a Real-Chess \texorpdfstring{$1/8$}{1/8}}\label{sec:real-eighth}

The $1/8$ position of Section~\ref{sec:targets} has $10$ men. Its e.p.-aware tree has $302$ end-position sides
with at most seven men, all resolved by the tablebase: in $290$ only one side has a passed pawn, of which $283$
agree with the model, $5$ disagree and $2$ are draws. A search for the larger end positions that decide the value
leaves, after one eight-man result against the model, one eight-man position and three of nine men; our
king-and-pawn solver (queen promotions only) shows that all four agree with the model on both sides. The hardest,
black to move in \pos{W b2 b5 / B a5 b4}, needed $6.1\times10^{10}$ states in $15.6$~h (Section~S7).

\begin{proposition}[A tablebase-backed oracle value of $1/8$; \lab{C1}]\label{prop:eighth-real}
For the $1/8$ position, with queen promotions only in our solver, $O=1/8$, whatever the open end positions do.
\end{proposition}

It assumes frozen kings; the full-chess value is open. For
$\upsec\star$ and $\tinyg{1}$ the lower bounds already equal the target values, with $10$ and $12$ end positions
of $9$ or $10$ men still open \lab{C1}. For positions A and B ($11$ men beside T) the deciding end positions with
$8$ men agree with the model \lab{C2}, but those with $9$ to $11$ men are unsolved, so their $\nim{2}$ is a model
value.
